% !TEX program = lualatex
% ---------------------------------------------------------------------------------------------
%  A guide to the architecture of the proof of RH announced by L. Alpöge  (riemann-rev)
%
%  Build:  python3 -B scripts/build_guide.py     (writes data.tex, runs LuaLaTeX, checks the layout)
%  The prose is hand-written. Numbers from the evidence, the status time and every link come from
%  data.tex. Each mathematical statement about the archive was checked against the Lean source
%  (pinned at commit 40e56bab); the declarations are named in the "Lean" line that closes a section.
%
%  Layout (reading edition): 11 pt Libertinus on a 140 mm measure, about 88 characters to the line;
%  formulas are displayed, enumerations are lists, section numbers hang in the margin. Only the
%  diagram is wider than the text. The page count is fixed in scripts/build_guide.py.
%
%  Rule of the text: every symbol and every name is explained where it first appears and listed
%  again in the notation on the last page; no letter has two meanings. Letters differ from the
%  archive's where the archive reuses one (its twist e is h here, its rank k is n_0, its Hecke
%  characters mu_j are xi_j, its family weights a are c).
% ---------------------------------------------------------------------------------------------
\documentclass[11pt,a4paper]{article}
\ifdefined\pdfvariable\pdfvariable suppressoptionalinfo \numexpr 1+2+4+8\relax\fi
\usepackage[a4paper,hmargin=35mm,top=22mm,bottom=27mm,footskip=12mm,nohead]{geometry}

\usepackage{fontspec}
\directlua{
  luaotfload.add_fallback("gfallback", {"STIXTwoMath-Regular.otf:mode=node;", "DejaVuSans.ttf:mode=node;"})
  luaotfload.add_fallback("gmonofallback", {"DejaVuSansMono.ttf:mode=node;", "STIXTwoMath-Regular.otf:mode=node;"})
}
\setmainfont{LibertinusSerif}[Extension=.otf,UprightFont=*-Regular,ItalicFont=*-Italic,
  BoldFont=*-Semibold,BoldItalicFont=*-SemiboldItalic,RawFeature={fallback=gfallback}]
\setsansfont{Inter}[Extension=.otf,UprightFont=*-Regular,ItalicFont=*-Italic,
  BoldFont=*-SemiBold,BoldItalicFont=*-SemiBoldItalic,RawFeature={fallback=gfallback}]
\setmonofont{JetBrainsMono}[Extension=.otf,UprightFont=*-Regular,ItalicFont=*-Italic,
  BoldFont=*-SemiBold,BoldItalicFont=*-SemiBoldItalic,Scale=MatchLowercase,
  RawFeature={-calt;-liga;fallback=gmonofallback}]
\usepackage{amsmath}
\usepackage[math-style=TeX]{unicode-math}
\setmathfont{LibertinusMath-Regular.otf}

\usepackage{xcolor}
\usepackage{tikz}
\usetikzlibrary{calc,arrows.meta,positioning,backgrounds}
\usepackage[protrusion=true,expansion=true]{microtype}
\usepackage{enumitem}
\usepackage{needspace}
\usepackage{booktabs}
\usepackage{array}
\usepackage{fancyhdr}
\usepackage{lastpage}
\usepackage{multicol}
\usepackage[hidelinks,pdfusetitle]{hyperref}

\definecolor{ink}{HTML}{0F172A}
\definecolor{inksoft}{HTML}{475569}
\definecolor{inkfaint}{HTML}{94A3B8}
\definecolor{rulegrey}{HTML}{CBD5E1}
\definecolor{arith}{HTML}{C2410C}
\definecolor{est}{HTML}{7C3AED}
\definecolor{ana}{HTML}{0891B2}
\definecolor{dest}{HTML}{E11D48}
\definecolor{ok}{HTML}{047857}
\definecolor{leancol}{HTML}{3730A3}

\title{The proof of the Riemann hypothesis announced by L. Alpöge: a guide to its architecture}
\author{riemann-rev}
\date{}

% generated by scripts/build_guide.py -- do not edit
\expandafter\newcommand\csname LN@main-rh\endcsname{\href{https://github.com/Deicyde/riemann-hypothesis/blob/40e56bab026552e089d7afcfeff08d2defb9c005/src/Main.lean\#L4-L4}{\leanfont riemannHypothesis}}
\expandafter\newcommand\csname LN@main-right\endcsname{\href{https://github.com/Deicyde/riemann-hypothesis/blob/40e56bab026552e089d7afcfeff08d2defb9c005/src/Main.lean\#L8-L8}{\leanfont LFunction\_ne\_zero\_of\_half\_lt\_re}}
\expandafter\newcommand\csname LN@main-strip\endcsname{\href{https://github.com/Deicyde/riemann-hypothesis/blob/40e56bab026552e089d7afcfeff08d2defb9c005/src/Main.lean\#L10-L10}{\leanfont LFunction\_eq\_zero\_re}}
\expandafter\newcommand\csname LN@imprimitive\endcsname{\href{https://github.com/Deicyde/riemann-hypothesis/blob/40e56bab026552e089d7afcfeff08d2defb9c005/src/Development/AutomorphicArithmetic/CharacterBoxes/Imprimitive.lean\#L5-L35}{\leanfont imprimitive}}
\expandafter\newcommand\csname LN@strip\endcsname{\href{https://github.com/Deicyde/riemann-hypothesis/blob/40e56bab026552e089d7afcfeff08d2defb9c005/src/Development/AutomorphicArithmetic/CharacterBoxes/ClaimStrip.lean\#L59-L73}{\leanfont Claim.strip}}
\expandafter\newcommand\csname LN@rh-right\endcsname{\href{https://github.com/Deicyde/riemann-hypothesis/blob/40e56bab026552e089d7afcfeff08d2defb9c005/src/Formalism/GammaFactorNonvanishing/RhOfRight.lean\#L5-L62}{\leanfont rh\_of\_right}}
\expandafter\newcommand\csname LN@G\endcsname{\href{https://github.com/Deicyde/riemann-hypothesis/blob/40e56bab026552e089d7afcfeff08d2defb9c005/src/Development/AutomorphicArithmetic/DedekindLFunctions/CharFamily.lean\#L10-L45}{\leanfont Deduction.G}}
\expandafter\newcommand\csname LN@mellin\endcsname{\href{https://github.com/Deicyde/riemann-hypothesis/blob/40e56bab026552e089d7afcfeff08d2defb9c005/src/Development/AutomorphicArithmetic/DedekindLFunctions/LandauTheorems/MellinRamp.lean\#L10-L32}{\leanfont MellinRamp}}
\expandafter\newcommand\csname LN@pole\endcsname{\href{https://github.com/Deicyde/riemann-hypothesis/blob/40e56bab026552e089d7afcfeff08d2defb9c005/src/Development/AutomorphicArithmetic/DedekindLFunctions/LandauTheorems/poleData.lean\#L248-L259}{\leanfont poleData}}
\expandafter\newcommand\csname LN@landau\endcsname{\href{https://github.com/Deicyde/riemann-hypothesis/blob/40e56bab026552e089d7afcfeff08d2defb9c005/src/Development/AutomorphicArithmetic/DedekindLFunctions/LandauTheorems/ZeroFreeOfRpowWindow.lean\#L11-L14}{\leanfont zeroFree\_of\_rpow\_window}}
\expandafter\newcommand\csname LN@window-consumer\endcsname{\href{https://github.com/Deicyde/riemann-hypothesis/blob/40e56bab026552e089d7afcfeff08d2defb9c005/src/Formalism/SharpRestriction/RamificationFrobenius/TotientGaloisClasses/TopWin.lean\#L68-L81}{\leanfont level\_of\_tiedWin}}
\expandafter\newcommand\csname LN@zero-top\endcsname{\href{https://github.com/Deicyde/riemann-hypothesis/blob/40e56bab026552e089d7afcfeff08d2defb9c005/src/Development/AutomorphicArithmetic/CharacterBoxes/ZeroFreeAllOfTopG.lean\#L26-L34}{\leanfont zeroFreeAll\_of\_topG}}
\expandafter\newcommand\csname LN@ischifield\endcsname{\href{https://github.com/Deicyde/riemann-hypothesis/blob/40e56bab026552e089d7afcfeff08d2defb9c005/src/Development/AutomorphicArithmetic/CharacterBoxes/ExistsPageN.lean\#L19-L25}{\leanfont IsChiField}}
\expandafter\newcommand\csname LN@chainint\endcsname{\href{https://github.com/Deicyde/riemann-hypothesis/blob/40e56bab026552e089d7afcfeff08d2defb9c005/src/Formalism/ThetaLiftIndependence/Statement.lean\#L40-L57}{\leanfont ChainInt}}
\expandafter\newcommand\csname LN@minor\endcsname{\href{https://github.com/Deicyde/riemann-hypothesis/blob/40e56bab026552e089d7afcfeff08d2defb9c005/src/Formalism/SharpRestriction/GrowthBounds/CoarseMultiplicity/KeyMinorDvd.lean\#L132-L155}{\leanfont key\_minor\_dvd}}
\expandafter\newcommand\csname LN@slope\endcsname{\href{https://github.com/Deicyde/riemann-hypothesis/blob/40e56bab026552e089d7afcfeff08d2defb9c005/src/Formalism/LatticeSlopes/LatticeSlopesProofs/SlopeInequality.lean\#L111-L125}{\leanfont slopeInequality\_of}}
\expandafter\newcommand\csname LN@chowla\endcsname{\href{https://github.com/Deicyde/riemann-hypothesis/blob/40e56bab026552e089d7afcfeff08d2defb9c005/src/Development/AutomorphicSpectrum/Proofs/GeometryProofs/ChowlaSelbergPeriodStatement.lean\#L9-L10}{\leanfont chowla\_selberg\_period}}
\expandafter\newcommand\csname LN@ask\endcsname{\href{https://github.com/Deicyde/riemann-hypothesis/blob/40e56bab026552e089d7afcfeff08d2defb9c005/src/Formalism/SizeWindowBrackets/FirstPrimeNeSecondPrime.lean\#L72-L92}{\leanfont ASK, MAIN}}
\expandafter\newcommand\csname LN@rank\endcsname{\href{https://github.com/Deicyde/riemann-hypothesis/blob/40e56bab026552e089d7afcfeff08d2defb9c005/src/Development/AbelianVarietyTheory/Definitions/GlobalTwisting/PairRowIffQdRows.lean\#L28-L32}{\leanfont RankVal}}
\expandafter\newcommand\csname LN@topg\endcsname{\href{https://github.com/Deicyde/riemann-hypothesis/blob/40e56bab026552e089d7afcfeff08d2defb9c005/src/Development/AutomorphicArithmetic/CharacterBoxes/ZeroFreeAllOfTopG.lean\#L12-L14}{\leanfont TopG}}
\expandafter\newcommand\csname LN@chbalance\endcsname{\href{https://github.com/Deicyde/riemann-hypothesis/blob/40e56bab026552e089d7afcfeff08d2defb9c005/src/Development/AutomorphicArithmetic/DedekindLFunctions/LandauTheorems/ChBalance.lean\#L200-L222}{\leanfont ch\_balance}}
\expandafter\newcommand\csname LN@shares\endcsname{\href{https://github.com/Deicyde/riemann-hypothesis/blob/40e56bab026552e089d7afcfeff08d2defb9c005/src/Formalism/SharpRestriction/RamificationFrobenius/TotientGaloisClasses/ExplicitFormulaErrorOfHostConditions.lean\#L133-L141}{\leanfont Frame.Shares}}
\expandafter\newcommand\csname LN@bounded\endcsname{\href{https://github.com/Deicyde/riemann-hypothesis/blob/40e56bab026552e089d7afcfeff08d2defb9c005/src/Formalism/SharpRestriction/RamificationFrobenius/TotientGaloisClasses/Analytic.lean\#L136-L156}{\leanfont BoundedCharsPointwise}}
\expandafter\newcommand\csname LN@pairex\endcsname{\href{https://github.com/Deicyde/riemann-hypothesis/blob/40e56bab026552e089d7afcfeff08d2defb9c005/src/Development/AutomorphicArithmetic/MeanSquareBounds/pairEx.lean\#L112-L169}{\leanfont QH.pairEx}}
\expandafter\newcommand\csname LN@core-gm\endcsname{\href{https://github.com/Deicyde/riemann-hypothesis/blob/40e56bab026552e089d7afcfeff08d2defb9c005/src/Development/AutomorphicArithmetic/PrimitiveCosets/CoreGm.lean\#L12-L74}{\leanfont core\_gm}}
\expandafter\newcommand\csname LN@supply\endcsname{\href{https://github.com/Deicyde/riemann-hypothesis/blob/40e56bab026552e089d7afcfeff08d2defb9c005/src/Development/AutomorphicArithmetic/CharacterBoxes/SupplyChi.lean\#L8-L16}{\leanfont supplyChi}}
\expandafter\newcommand\csname LN@vwin-def\endcsname{\href{https://github.com/Deicyde/riemann-hypothesis/blob/40e56bab026552e089d7afcfeff08d2defb9c005/src/Development/AutomorphicArithmetic/CharacterBoxes/ExistsPageN.lean\#L44-L47}{\leanfont VWinAt}}
\expandafter\newcommand\csname LN@window1\endcsname{\href{https://github.com/Deicyde/riemann-hypothesis/blob/40e56bab026552e089d7afcfeff08d2defb9c005/src/Development/AutomorphicArithmetic/AutomorphicVarieties/LayeredIsomorphisms/Window1.lean\#L6-L58}{\leanfont window1}}
\expandafter\newcommand\csname LN@behrend\endcsname{\href{https://github.com/Deicyde/riemann-hypothesis/blob/40e56bab026552e089d7afcfeff08d2defb9c005/src/Formalism/FrobeniusOrbits/FrobeniusOrbitsProofs/ThickExists.lean\#L181-L187}{\leanfont hostThickExists}}
\expandafter\newcommand\csname LN@row-def\endcsname{\href{https://github.com/Deicyde/riemann-hypothesis/blob/40e56bab026552e089d7afcfeff08d2defb9c005/src/Development/AutomorphicArithmetic/Foundations/GaloisRepresentations/TypeConditions/WeightsConjBis.lean\#L69-L73}{\leanfont Row0b}}
\expandafter\newcommand\csname LN@row-prod\endcsname{\href{https://github.com/Deicyde/riemann-hypothesis/blob/40e56bab026552e089d7afcfeff08d2defb9c005/src/Development/AutomorphicArithmetic/Foundations/GaloisRepresentations/IntegralStructures/RowSrc.lean\#L6-L6}{\leanfont galoisLocalGlobalStrong\_recParameter}}
\expandafter\newcommand\csname LN@seven\endcsname{\href{https://github.com/Deicyde/riemann-hypothesis/blob/40e56bab026552e089d7afcfeff08d2defb9c005/src/Development/AutomorphicArithmetic/ThetaCorrespondence/ThetaPairLifts/ThirdPlaceSeries/EveryCharacterOfLiteratureOfVWin.lean\#L810-L813}{\leanfont every\_character\_of\_literature\_of\_vWin}}
\newcommand{\FullProofIcon}{\textcolor{ok}{\checkmark\checkmark}}
\newcommand{\FullProofLine}{two independent kernels accepted the complete proof. In runs of the reading project the Rust kernel NanoDA accepted all 1{,}484{,}733 declarations of the exported proof, and the five statements matched an independently written Mathlib challenge. S.\,DeDeo and A.\,Sutherland each report a Lean kernel replay of all 175{,}187 modules. No axioms beyond {\leanfont propext}, {\leanfont Classical.choice}, {\leanfont Quot.sound}, the three standard axioms of Lean.}
\newcommand{\Rmax}{0.023}
\newcommand{\RHband}{0.0462}
\newcommand{\ModuleCount}{175{,}187}
\newcommand{\DirichletModules}{1{,}421}
\newcommand{\NanodaDecls}{61{,}697}
\newcommand{\AsOf}{20 November 2026, 23:51 UTC}


% ------------------------------------------------------------------------------------- page
\linespread{1.08}                       % 11 pt on 14.7 pt
\setlength{\parindent}{0pt}
\setlength{\parskip}{5.5pt plus 1.5pt minus 0.5pt}
\setlength{\abovedisplayskip}{7.5pt plus 2pt minus 3pt}
\setlength{\belowdisplayskip}{7.5pt plus 2pt minus 3pt}
\setlength{\abovedisplayshortskip}{3pt plus 2pt}
\setlength{\belowdisplayshortskip}{6pt plus 2pt minus 2pt}
\emergencystretch=1.5em
\clubpenalty=10000 \widowpenalty=10000 \displaywidowpenalty=10000
\relpenalty=7000 \binoppenalty=9000     % keep short formulas on one line
\raggedbottom
\pagestyle{fancy}\fancyhf{}\renewcommand{\headrulewidth}{0pt}
\fancyfoot[C]{\sffamily\fontsize{7.6}{9}\selectfont\color{inkfaint}\thepage\,/\,\pageref*{LastPage}}

\setlist{topsep=2pt,partopsep=0pt,parsep=0pt,itemsep=5pt,leftmargin=2.1em,labelsep=0.55em}

% ------------------------------------------------------------------------------------- macros
\newcommand{\leanfont}{\ttfamily\fontsize{8.6}{10.5}\selectfont\color{leancol}}
\newcommand{\lean}[1]{\csname LN@#1\endcsname}
\newcommand{\Sha}{\text{\normalfont Ш}}
\newcommand{\eps}{\varepsilon}
\AtBeginDocument{\renewcommand{\Re}{\operatorname{Re}}}
\newcommand{\Q}{\mathbb{Q}}
\newcommand{\Z}{\mathbb{Z}}
\newcommand{\F}{\mathbb{F}}
\newcommand{\LL}[1]{\frac{L'}{L}(#1)}
\newcommand{\chid}{\chi_{-d}}
\newcommand{\chiq}{\chi_{-q}}
\newcommand{\len}{\operatorname{len}}
\newcommand{\Gal}{\operatorname{Gal}}
\newcommand{\kw}[1]{\emph{#1}}
\newcommand{\m}[1]{\mbox{$#1$}}        % a short formula that must not be broken across lines
\newcommand{\keep}{\par\nopagebreak}    % new paragraph, kept on the page of the previous one
% a section number in a coloured disc: #1 number, #2 colour, #3 diameter, #4 font size
\newcommand{\disc}[4]{\tikz[baseline=(n.base)]\node[circle,fill=#2,inner sep=0pt,minimum size=#3,
  text=white,font=\sffamily\bfseries\fontsize{#4}{#4}\selectfont](n){#1};}
\newcommand{\dotnum}[2]{\disc{#1}{#2}{12.4pt}{7.6}}
% numbered section: the disc hangs in the left margin, the title starts at the text edge
\newcommand{\sect}[3]{\par\addvspace{17pt plus 4pt minus 2pt}\Needspace{6\baselineskip}%
  \noindent\llap{\dotnum{#1}{#2}\hspace{7pt}}{\fontsize{13.2}{16}\selectfont\bfseries\color{ink}#3}\par\nopagebreak\vspace{-1pt}}
\newcommand{\head}[1]{\par\addvspace{17pt plus 4pt minus 2pt}\Needspace{6\baselineskip}%
  \noindent{\fontsize{13.2}{16}\selectfont\bfseries\color{ink}#1}\par\nopagebreak\vspace{-1pt}}
% the line that closes a section: where its statements live in the Lean archive
\newcommand{\leantag}{{\sffamily\bfseries\fontsize{6.3}{8}\selectfont\color{leancol}\addfontfeatures{LetterSpace=7}LEAN}}
\newcommand{\sep}{\unskip\nobreak\hspace{5pt}{\color{inkfaint}\textperiodcentered}\hspace{5pt plus 2pt}\ignorespaces}
\newcommand{\leanline}[1]{\par\nopagebreak\vspace{1pt}{\fontsize{8.9}{12.4}\selectfont\color{inksoft}\raggedright
  \noindent\llap{\leantag\hspace{7pt}}#1\par}}
\newcommand{\leanin}[1]{\par\vspace{3.5pt}{\fontsize{8.9}{12.4}\selectfont\color{inksoft}\raggedright
  \noindent\leantag\hspace{7pt}#1\par}}
% a displayed statement in a tinted box: #1 colour, #2 text
\newcommand{\statement}[2]{\par\addvspace{7pt}\begingroup\setlength{\fboxsep}{8.5pt}\setlength{\fboxrule}{0.6pt}%
  \noindent\fcolorbox{#1!30!white}{#1!3!white}{\parbox{\dimexpr\linewidth-2\fboxsep-2\fboxrule\relax}{#2}}%
  \endgroup\par\addvspace{7pt}}
% diagram: title line of a box (disc with the section number, coloured sans-serif title)
\newcommand{\bt}[3]{{\sffamily\bfseries\fontsize{7.4}{9.4}\selectfont\color{#2!82!black}%
  \mbox{\disc{#1}{#2}{10.4pt}{6.6}\hspace{2.6pt}#3}}\\[2pt]}
\newcommand{\nt}[2]{\par\noindent\hangindent=1.1em\hangafter=1 #1\enspace{\color{inksoft}#2}\par\vspace{1pt}}
\newcommand{\legdot}[1]{\tikz[baseline=-0.62ex]\fill[#1] (0,0) circle (2.5pt);\hspace{3pt}}

\begin{document}
\color{ink}

% ------------------------------------------------------------------------------------- title
{\raggedright
{\fontsize{21.5}{25.5}\selectfont\bfseries The proof of the Riemann hypothesis\\ announced by L.\,Alpöge\par}
\vspace{5pt}
{\fontsize{12.4}{16}\selectfont\itshape\color{inksoft} Its architecture: the statements, the flow, and why each idea calls for
the next\par}}
\vspace{8pt}
{\color{rulegrey}\hrule height 0.6pt}
\vspace{7pt}
\thispagestyle{empty}\enlargethispage{12mm}

\textbf{One hard estimate and two classical transformations of it.} Arithmetic geometry bounds a smoothed sum over
primes \emph{from above}, by a square root, on sparse but long intervals. A localized form of Landau's theorem
turns this one-sided, sparse information into a zero-free half-plane, and functional equations carry the half-plane to
the critical line. The text reads the proof backwards from the theorem; the diagram shows the same flow forwards.

% ------------------------------------------------------------------------------------- diagram
\par\vspace{6pt}
\noindent\makebox[\linewidth][c]{%
\begin{tikzpicture}[x=1cm,y=1cm,
  box/.style 2 args={draw=#1!55!white,line width=0.6pt,rounded corners=3pt,fill=#1!4!white,
    minimum height=1.8cm,text width=#2,inner xsep=4.5pt,inner ysep=4pt,align=flush left,
    font=\fontsize{7.9}{9.6}\selectfont},
  low/.style={minimum height=1.52cm},
  flow/.style={-{Stealth[length=4.8pt,width=4.1pt]},line width=0.8pt,draw=ink!78},
  lab/.style={font=\sffamily\fontsize{6.9}{8}\selectfont,text=ink!85,inner sep=1.5pt,fill=white},
  cut/.style={dash pattern=on 2.8pt off 2.2pt,line width=0.65pt,draw=ok},
  cutlab/.style={font=\sffamily\fontsize{6.9}{8.6}\selectfont,text=ok,inner sep=1pt}]
  % upper row: four stations (3.52 cm wide, 1.18 cm apart)
  \node[box={arith}{3.2cm}] (B1) at (1.775,0)  {\bt{7}{arith}{Automorphic input}Galois representations
     $\rho_\pi$ over CM fields; Siegel--Weil formula; local Langlands};
  \node[box={est}{3.2cm}]   (B2) at (6.475,0)  {\bt{4}{est}{Chain of theta lifts}$U(X_0)\to\cdots\to U(X_m)$,
     $m\asymp\sqrt x$; divisibility; Bost's slope inequality};
  \node[box={est}{3.2cm}]   (B3) at (11.175,0) {\bt{5}{est}{Character averaging}Fourier analysis on
     $\operatorname{Gal}(E/\Q)$; positive weights; twists discarded};
  \node[box={ana}{3.2cm}]   (B4) at (15.875,0) {\bt{2}{ana}{\hspace{-1pt}\disc{3}{ana}{10.4pt}{6.6}\hspace{2.6pt}Localized Landau}positive
     residues; high derivatives on windows; $\theta$ after the zero, $\eps$ after $\theta$};
  % lower row
  \node[box={arith}{7.9cm},low] (C1) at (4.125,-2.52) {\bt{6}{arith}{Iwasawa theory}total length of the groups $\Sha^1_\ell(E,h)$ at most
     $\exp(x^{\eps}/2)$ on a whole window $[Y,Y^2]$; Ferrero--Washington; interval lemma};
  \node[box={est}{3.2cm},low]   (C2) at (11.175,-2.52) {\bt{5}{est}{Hosts}CM fields $E$ in a cyclotomic
     tower over $K_\psi(\sqrt{-q},\sqrt{-d})$; $q$, $d$ by mean squares};
  \node[box={dest}{3.2cm},low]  (C3) at (15.875,-2.52) {\bt{1}{dest}{Reflection}Euler factors, functional
     equations: RH and GRH for Dirichlet $L$-functions};
  % the two verified cuts
  \draw[cut] (8.825,1.08) -- (8.825,-3.5);
  \draw[cut] (13.525,1.08) -- (13.525,-3.5);
  \node[cutlab,anchor=north east,align=right] at (8.73,-3.44)
    {\textbf{\checkmark\ Lean kernel:}\ \ \texttt{TopG} $\wedge$ $\forall\psi$ \texttt{VWinAt}\,$\psi$\,401 $\Longrightarrow$ Theorem};
  \node[cutlab,anchor=north west,align=left] at (13.62,-3.44)
    {\textbf{\checkmark\checkmark\ Lean and an independent}\\\textbf{Rust kernel:}\ \ $(\star)$ for $\zeta$ $\Longrightarrow$ RH};
  % arrows
  \draw[flow] (B1.east) -- (B2.west) node[midway,above=2.6pt,lab]{periods};
  \draw[flow] (B2.east) -- (B3.west) node[midway,above=2.6pt,lab]{\texttt{TopG}};
  \draw[flow] (B3.east) -- (B4.west) node[midway,above=2.6pt,lab]{$(\star)$};
  \draw[flow] (B4.south) -- (C3.north) node[midway,right=2pt,lab]{$\Re s>\frac12$};
  \draw[flow] (B1.south) -- (B1.south |- C1.north) node[midway,right=2pt,lab]{$\rho_\pi$, base change};
  \draw[flow] (B2.south |- C1.north) -- (B2.south) node[midway,right=2pt,lab]{loss of divisibility};
  \draw[flow] (C1.east) -- (C2.west) node[midway,above=2.6pt,lab]{\texttt{VWinAt}};
  \draw[flow] (C2.north) -- (B3.south) node[midway,right=2pt,lab]{hosts};
  % legend of the four strata
  \node[anchor=south west,inner sep=0pt,font=\sffamily\fontsize{6.9}{8}\selectfont,text=inksoft] at (0.02,1.2)
    {\legdot{arith}arithmetic supply\hspace{11pt}\legdot{est}the estimate\hspace{11pt}\legdot{ana}the
     analytic engine\hspace{11pt}\legdot{dest}the destination};
\end{tikzpicture}}\par
\vspace{4pt}
{\fontsize{9}{12}\selectfont\itshape\color{inksoft} The flow of the argument; circled numbers refer to the sections
below. Everything to the right of a dashed line has been machine-checked as an implication from the statements that
cross the line.\par}

% ------------------------------------------------------------------------------------- theorem
\statement{dest}{\textbf{Theorem.}\enspace\emph{Let $\chi$ be a Dirichlet character of modulus $N\ge1$. Then
$L(s,\chi)\neq0$ for $\Re s>\frac12$, except at the pole $s=1$ of a principal character. Consequently every zero of
$L(s,\chi)$ with $0<\Re s<1$ lies on the critical line $\Re s=\frac12$; for $N=1$, where $L(s,\chi)=\zeta(s)$, this is
the Riemann hypothesis.}%
\leanin{\lean{main-right}\sep\lean{main-strip}\sep\lean{main-rh}, a term of Mathlib's {\leanfont RiemannHypothesis}}}
{\fontsize{8.9}{12.2}\selectfont\color{inksoft}\noindent Lines marked \leantag\ link to the Lean archive, the formal
proof; Lean and all notation are explained at the end.\par}

\head{The shape of the argument: Stepanov over $\Z$}
RH is square-root cancellation among the primes. For curves over a finite field~$\F$, Stepanov obtained it from an
auxiliary function vanishing to high order at the rational points: comparing zeros with the degree bounds the number
of points from \emph{above}, and upper bounds over all finite extensions of~$\F$ suffice.

The announcement calls its argument ``a descendant of Stepanov's'', each theta lift ``\mbox{mimicking} multiplicity'',
with $\sqrt X$ lifts for the von Mangoldt sum of length~$X$.

\begingroup\par\nopagebreak\vspace{-2pt}\fontsize{10.4}{13.8}\selectfont
\noindent\begin{tabular}{@{}>{\raggedright\arraybackslash}p{0.34\linewidth}@{\hspace{0.025\linewidth}}c@{\hspace{0.025\linewidth}}>{\raggedright\arraybackslash}p{0.565\linewidth}@{}}
\toprule
{\sffamily\fontsize{7.6}{10}\selectfont\color{inksoft}STEPANOV, FOR CURVES OVER $\F$} & &
{\sffamily\fontsize{7.6}{10}\selectfont\color{inksoft}THE ANNOUNCED PROOF, OVER $\Z$}\\
\midrule
vanishing to order $m$ & $\longrightarrow$ & divisibility of an $m$-fold iterated theta lift (§4)\\
counting degrees & $\longrightarrow$ & Bost's slope inequality (§4)\\
the finite extensions of $\F$ & $\longrightarrow$ & the layers of a cyclotomic tower of fields (§4)\\
``upper bounds suffice'' & $\longrightarrow$ & a theorem of Landau type (§3)\\
\bottomrule
\end{tabular}\par\endgroup

% ------------------------------------------------------------------------------------- 1
\sect{1}{dest}{Reflection}
Let $\chi$ be a Dirichlet character modulo~$N$ and $\chi^*$ the primitive character that induces it. Then
\[
L(s,\chi)=L(s,\chi^*)\prod_{p\mid N}\bigl(1-\chi^*(p)\,p^{-s}\bigr),
\]
the product running over the primes~$p$ dividing~$N$, and no factor of the product vanishes for $\Re s>0$.

For primitive~$\chi$ let $\Lambda(s,\chi)$ be the completed $L$-function, that is, $L(s,\chi)$ multiplied by its Gamma
factor and a power of the modulus, and let $W(\chi)$ be the root number, a complex number of absolute value~$1$. With
$\bar\chi$ the complex conjugate character, the functional equation
\[
\Lambda(s,\chi)=W(\chi)\,\Lambda(1-s,\bar\chi)
\]
sends a zero in the critical strip \m{0<\Re s<1} with \m{\Re s<\frac12} to a zero of $L(\,\cdot\,,\bar\chi)$ with
\m{\Re s>\frac12}.
\keep
So it suffices to prove \m{L(s,\psi)\neq0} for \m{\Re s>\frac12} and every primitive character~$\psi$, the trivial
one, for which $L=\zeta$, included.
\leanline{all characters \lean{imprimitive}\sep reflection \lean{strip}\sep reflection for $\zeta$ \lean{rh-right}}

% ------------------------------------------------------------------------------------- 2
\sect{2}{ana}{The statistic}
Let $\Lambda(n)$ be the von Mangoldt function, equal to $\log p$ if $n$ is a power of the prime~$p$ and to $0$
otherwise, and let $\gamma$ be Euler's constant. For a real number \m{y\ge1} and a primitive character \m{\chi\ne1} put
\begin{align*}
G(y)&=\sum_{n\le y}\Lambda(n)\Bigl(\frac yn-1\Bigr)-y\,(\log y-1-\gamma),\\
g_\chi(y)&=\sum_{n\le y}\Lambda(n)\,\chi(n)\Bigl(\frac yn-1\Bigr)+y\,\LL{1,\chi}.
\end{align*}
The weight $(y/n-1)_+$, where \m{(u)_+:=\max(u,0)}, is a \kw{ramp}: it vanishes at \m{n=y} and grows as $n$ decreases.
The last term of each line is the \kw{centring}; it is what makes the transform below regular at \m{s=1}.

The \kw{Mellin transform} of a function~$f$ on $[1,\infty)$ is the function
\m{s\mapsto\int_1^\infty f(y)\,y^{-s-1}\,dy}. That of the ramp $(y/n-1)_+$ is $n^{-s}/(s(s-1))$, so summation
over~$n$ gives, for \m{\Re s>1},
\[
\int_1^\infty g_\chi(y)\,y^{-s-1}\,dy=-\frac{1}{s(s-1)}\Bigl[\LL{s,\chi}-s\,\LL{1,\chi}\Bigr],
\]
and a similar formula for~$G$, with $\zeta$ in place of $L(\,\cdot\,,\chi)$. These transforms continue to \m{\Re s>0},
where their only singularities are simple poles at the zeros~$\rho$ of the $L$-function: if $\rho$ is a zero of
multiplicity $m_\chi(\rho)$, the residue is $-m_\chi(\rho)/(\rho(\rho-1))$. Read backwards, this is the explicit formula
\[
G(y)=\log2\pi-\sum_\rho\frac{y^\rho}{\rho(\rho-1)}+O(y^{-2}),
\]
in which $\rho$ runs over the zeros of~$\zeta$ in the critical strip, with multiplicity. Hence RH gives
\m{\lvert G(y)-\log2\pi\rvert\le\RHband\,\sqrt y+O(y^{-2})}, the constant being \m{2+\gamma-\log4\pi}. Numerically one
finds \m{\lvert G(y)-\log2\pi\rvert<\Rmax\sqrt y} for \m{10\le y\le10^7}.
\keep
The proof needs far less: an upper bound, on windows.
\leanline{the two sums \lean{G}\sep transform of the ramp \lean{mellin}\sep poles and residues \lean{pole}}

% ------------------------------------------------------------------------------------- 3
\sect{3}{ana}{A localized Landau theorem}
Landau's theorem says that the Mellin transform of a function of constant sign has a singularity at the real point
of its abscissa of convergence. Applied to \m{Cy^{\theta}-G(y)} it shows that a bound \m{G(y)\le Cy^{\theta}}
\emph{for all large~$y$} would exclude zeros of~$\zeta$ with \m{\Re s>\theta}. The arithmetic of §4--6 yields such a
bound only on certain intervals, and only for several characters together. So the theorem has to be localized, and
the zeros of different $L$-functions must be kept from cancelling.

A \kw{positive family} $F$ is a finite set of Dirichlet characters, each either non-principal or the trivial character
of modulus one, with weights \m{c_\chi>0}. Put
\[
G_F(y)=\sum_{\chi\in F}c_\chi\,\Re g_\chi(y),\qquad\text{with }g_1:=G\text{ for the trivial character.}
\]
A \kw{window} is an interval $[Y,Y^{1+\delta}]$ with \m{\delta>0}.

\statement{ana}{\textbf{Proposition.}\enspace\emph{Let $F$ be a positive family, $0<\theta<1$ and $\delta>0$. Suppose
that some $L(s,\chi)$ with $\chi\in F$ has a zero with $\Re s>\theta$. Then for every real number~$C$ there is $Y_0$
such that every window $[Y,Y^{1+\delta}]$ with $Y\ge Y_0$ contains a point~$y$ with $G_F(y)>C\,y^{\theta}$.}}

Thus a zero shows itself, with a definite sign, in every late window. Three ideas make this work.
\begin{enumerate}[label=\textup{(\roman*)}]
\item \kw{Zeros cannot cancel.} By §2, and because \m{\Re g_\chi=\frac12(g_\chi+g_{\bar\chi})}, the transform of
$G_F$ has at a zero~$\rho$ the residue $-S(\rho)/(\rho(\rho-1))$, where
\[
S(\rho)=\tfrac12\sum_{\chi\in F}c_\chi\bigl(m_\chi(\rho)+m_{\bar\chi}(\rho)\bigr)>0.
\]
As all weights are positive, the contributions of the different characters add up. At a real zero the residue is
positive, which covers possible real zeros close to~$1$ (Siegel zeros).
\item \kw{High derivatives localize.} Put \m{g(u)=G_F(e^u)} and let \m{\widehat g(s)=\int_0^\infty g(u)\,e^{-su}\,du}
be its Laplace transform; this is the transform of §2 in the variable \m{u=\log y}. For an integer \m{k\ge0},
\[
A_k(s):=\frac{(-1)^k}{k!}\,\widehat g^{\,(k)}(s)=\int_0^\infty g(u)\,\frac{u^k}{k!}\,e^{-su}\,du.
\]
Write \m{\beta=\Re s}. The weight $u^ke^{-\beta u}/k!$ is a bump centred at \m{u=k/\beta}, of relative width
$k^{-1/2}$. Taking \m{k=\lceil(1+\delta/2)\,\beta\log Y\rceil} places it on the window, where
\m{\log Y\le u\le(1+\delta)\log Y}.
\item \kw{One side suffices.} With \m{g^+=\max(g,0)} one has \m{\lvert g\rvert+g=2g^+}. So an upper bound
\m{g(u)\le Ce^{\theta u}} on the window gives, for every \m{\eta>0} and a constant $C_\eta$,
\[
\lvert A_k(s)\rvert+A_k(\beta)\ \le\ C_\eta\,Y^{\eta}(\beta-\theta)^{-k}.
\]
On the other hand a pole of $\widehat g$ closer to~$s$ than \m{\beta-\theta} makes the left side grow faster in~$k$
than $(\beta-\theta)^{-k}$: through the first term if the pole is not real, and through the second, with the right sign
by~(i), if it is. For large~$k$ the two statements are incompatible.
\end{enumerate}

\vspace{3pt}
\textbf{The order of choices.} Let $\psi$ be a primitive character of order~$o$, that is, \m{\psi^o=1} and no smaller
power is trivial. Let $P(\psi)$ be the positive family formed by the primitive characters inducing the $o$ powers
of~$\psi$, with weight~$1$ on the trivial character and $1/(3o+1)$ on each of the others; write
\m{G_\psi=G_{P(\psi)}}. The reason for these weights is in §5. The arithmetic half of the proof is:

\statement{est}{\textbf{Statement $(\star)$.}\enspace\emph{There is $\eps_1>0$ such that for every $0<\eps<\eps_1$
there are arbitrarily large~$Y$ with}\par\vspace{5pt}
\centerline{$G_\psi(x)\le x^{1/2+9\eps}$\quad\emph{for all}\quad$x\in[Y,Y^{3/2}]$.}}

Suppose $L(s,\psi)$ had a zero~$\rho$ with \m{\Re\rho>\frac12}. Choose $\theta$ with \m{\frac12<\theta<\Re\rho}, and
then $\eps$ with \m{\frac12+9\eps\le\theta}. By $(\star)$ there are arbitrarily late windows $[Y,Y^{3/2}]$ on which
\m{G_\psi(x)\le x^{\theta}}. This contradicts the Proposition for the family $P(\psi)$, with \m{\delta=\frac12} and
\m{C=1}.
\keep
So $\theta$ is chosen after the zero and $\eps$ after~$\theta$. No bound at \m{\eps=0} and no uniformity in~$\eps$ is
needed, and the windows may depend on everything.
\leanline{the Proposition \lean{landau}\sep the order of choices \lean{window-consumer}, \lean{zero-top}}

% ------------------------------------------------------------------------------------- 4
\sect{4}{est}{The estimate, I: a slope inequality for a chain of theta lifts}
\textbf{The principle, in its simplest case.} Let $K$ be an imaginary quadratic field and $\vartheta$ a theta function
on an elliptic curve with complex multiplication by~$K$, with period~$\Omega$ and invariant derivation~$D$. The
announcement records (entry~2) that for every prime~$p$ inert in~$K$ and every integer \m{a\ge1} the number
$D^a\vartheta(0)/\Omega^a$ is divisible by
\[
p^{\,a/p+a/p^3+\cdots-O(\log a)},
\]
much as $a!$ is divisible by $p^{\lfloor a/p\rfloor+\lfloor a/p^2\rfloor+\cdots}$. Now a non-zero algebraic integer
that is divisible by many primes cannot be small: by the product formula, the sum of $\log p$ times these exponents,
over the inert primes, is at most the logarithm of its size at the infinite places. The left side of this
inequality is a sum over primes weighted by a ramp, as in §2; the right side can be computed. So one obtains an
\emph{upper} bound for a prime sum. The rest of this section makes the error terms small enough: the single
derivative becomes a chain of $m$ theta lifts, and the product formula becomes a slope inequality.

\textbf{The setting.} Fix a primitive character~$\psi$ of order~$o$, a small \m{\eps>0} and a \kw{scale}~$x$, a large
real number: the aim is $(\star)$ at the point~$x$. Choose two auxiliary primes $q$ and~$d$, subject to conditions
given in §5, and put \m{K=\Q(\sqrt{-d})}. Let $\chid$ be the quadratic Dirichlet character of~$K$: \m{\chid(p)=1}
if the prime~$p$ splits in~$K$, $-1$ if it is inert, and $0$ if \m{p=d}. Define $\chiq$ likewise for
$\Q(\sqrt{-q})$.

\textbf{The host field.} Let $e$ be an integer with \m{x^{2+8\eps}\le2^e\le2x^{2+8\eps}}. The \kw{host} is the abelian
number field
\[
E=\Q(\zeta_{2^{e+2}})^{+}\cdot K_\psi\cdot\Q(\sqrt{-q},\sqrt{-d}).
\]
Here $\zeta_n$ is a primitive $n$-th root of unity, the superscript~$+$ denotes the maximal real subfield, and
$K_\psi$ is the field cut out by~$\psi$, that is, the fixed field of the kernel of~$\psi$; it is cyclic of degree~$o$
over~$\Q$. The field~$E$ is a CM field: a totally imaginary quadratic extension of a totally real field~$E^+$, with
complex conjugation~$c$. As $x$ grows, the first factor climbs the tower
\m{\Q\subset\Q(\zeta_8)^+\subset\Q(\zeta_{16})^+\subset\cdots}. This cyclotomic tower plays the part of the finite
extensions of~$\F$ in Stepanov's proof.

\textbf{The chain.} Let \m{X_0,X_1,\dots,X_m} be Hermitian spaces over~$E$, for the conjugation~$c$, that are positive
definite at every real place of~$E^+$ (``totally definite'' in the announcement), with
\[
\dim X_j=n_0+j\,(t+1),\qquad n_0\asymp x^{9/2+10\eps}\ \text{even},\qquad t\ \text{odd},\ \ n_0<t\le x^{400}.
\]
Let $U(X_j)$ be the unitary group of~$X_j$. Because $X_j$ is definite, $U(X_j)$ is compact at the infinite places, and
an automorphic form on it, of fixed level and weight, is simply a function~$f$ on a finite set: the classes~$M$ of
Hermitian lattices in a genus, each counted with the weight $1/\lvert\operatorname{Aut}M\rvert$.

The \kw{theta correspondence} for the pair $U(X_j)$, $U(X_{j+1})$, together with a Hecke character $\xi_j$ of~$E$,
provides a linear map from forms on $U(X_j)$ to forms on $U(X_{j+1})$, the \kw{theta lift}
\[
(\Theta_jf)(M')=\sum_{M}\frac{\vartheta_j(M,M')}{\lvert\operatorname{Aut}M\rvert}\,f(M),
\]
where $M$ and $M'$ run over the classes of lattices in $X_j$ and $X_{j+1}$, and $\vartheta_j$ is a theta kernel
twisted by~$\xi_j$. The \kw{chain} is the composite \m{\Theta=\Theta_{m-1}\circ\cdots\circ\Theta_0}. In terms of
automorphic representations each lift adjoins to a representation~$\pi$ a block attached to~$\xi_j$; the announcement
writes \m{\pi\mapsto\pi\boxplus\mu_j[t]} (entry~8, with $\mu_j$ for our~$\xi_j$), where $\boxplus$ is the isobaric sum.
The number~$m$ of lifts is the parameter to be optimized.

Three computations enter the slope inequality.
\begin{enumerate}[label=\textit{(\alph*)}]
\item \kw{Finite places: divisibility.} The forms with integral values make up a lattice in each space of forms. The
theta kernels are integral and the factors $1/\lvert\operatorname{Aut}M\rvert$ have bounded denominators, so the
chain maps integral forms to integral forms after multiplication by an explicit integer. The point is that the image
is far more divisible than that, at the primes inert in~$K$, as in the simplest case above. The local source of this
divisibility is a statement on minors: if \m{\varpi^2=p} in a commutative ring, then $\varpi$ divides a certain
$2\times2$ minor formed from two rows of the lifted data.

Each prime being weighted by its logarithm, the divisibility gained at \kw{height} \m{a>0}, the analogue of the
order~$a$ of the derivative, is
\[
\mathrm{ASK}(a)=\sum_{n\ge1}\Lambda(n)\,\bigl(1-\chid(n)\bigr)\Bigl(\frac an-w_n\Bigr)_{\!+}.
\]
The factor $1-\chid(n)$ equals $2$ at the odd powers of primes inert in~$K$, $1$ at the powers of~$d$, and $0$ at all
other prime powers. The \kw{weights} are the values of a positive function~$w$ on $\Phi:=\Gal(E/K)$, and $w_n$ is its
value at an element determined by the Frobenius class of~$n$ (for \m{n=p^i}, the $i$-th power of the Frobenius
element of~$p$ in $\Gal(E/\Q)$). The weights that matter are close to a common value \m{b:=\exp(x^{\eps})}, and the heights that occur satisfy
\m{a/b\approx x}. Since \m{(a/n-w)_+=w\,(y/n-1)_+} for \m{y=a/w}, each term is $w$
times the ramp of §2 at a point~$y$ close to~$x$.
\item \kw{The product formula: slopes.} A \kw{Euclidean lattice}~$\mathcal L$ is a free abelian group of finite rank
with a positive definite quadratic form. Its \kw{slope} is
\m{\mu(\mathcal L)=-\log\operatorname{covol}(\mathcal L)/\operatorname{rk}(\mathcal L)}, and $\mu_{\max}(\mathcal L)$
is the largest slope of a non-zero sublattice. For an injective homomorphism
\m{\varphi\colon\mathcal L\to\mathcal L'} of Euclidean lattices, with operator norm $\lVert\varphi\rVert$, Bost's
slope inequality reads
\[
\mu(\mathcal L)\le\mu_{\max}(\mathcal L')+\log\lVert\varphi\rVert.
\]
It is applied to the lattices of integral forms and to the map between them given by the chain. What the map gains
at the finite places, through divisibility, it must pay at the infinite ones, through its norm.
\item \kw{Infinite places: periods.} The size of the chain at the infinite places is computed exactly. The
announcement names the tools (entries 9, 12 and~18): a seesaw identity with the Rallis inner product formula for the
norms of the lifts, Siegel's mass formula, Weyl's formula, and the Chowla--Selberg formula for the periods at the
points with complex multiplication. For an elliptic curve with complex multiplication by~$K$ and its suitably
normalized period~$\Omega$, the last one reads
\[
\log\Omega=\gamma+\log\frac{2\pi}{\sqrt d}-\LL{1,\chid}.
\]
Together with the growth \m{a\log a-a} of derivatives of order~$a$ (Stirling's formula), this produces the main
term
\[
\mathrm{MAIN}(a)=a\,\Bigl(\log a-1-\gamma+\LL{1,\chid}-\overline{\log w}\,\Bigr),
\]
where $\overline{\log w}$ is the mean value of $\log w$ over~$\Phi$.
\end{enumerate}

\textbf{The outcome.} Applied to the chain, divided by the common weight~$b$ and averaged over the $m$ lifts, the
slope inequality becomes
\[
\frac1b\,\Bigl(\mathrm{ASK}(a)-\mathrm{MAIN}(a)\Bigr)\ \le\ x^{O(\eps)}\Bigl(\frac xm+m+\sqrt x\Bigr).
\]
On the right, the errors of the whole chain are shared among its $m$ lifts, which gives $x/m$; each lift has a cost
that grows with the length of the chain, which gives~$m$; the third term does not depend on~$m$. A short chain gains
too little and a long one costs too much: the optimum is \m{m\asymp\sqrt x}. In the archive
\m{\sqrt{x\delta_1}\le m\le2\sqrt{x\delta_1}} with \m{\delta_1=x^{-\eps/8}}.

On the left, suppose for a moment that all weights are equal to one number~$w$, and put \m{y=a/w}. The definitions
of §2 then give at once
\[
\mathrm{ASK}(a)-\mathrm{MAIN}(a)=w\,\bigl(G(y)-g_{\chid}(y)\bigr).
\]
\emph{The centring of $G$ and of the $g_\chi$, Euler's constant included, is the archimedean side of a slope
inequality.} So the outcome is an upper bound of size about $\sqrt x$, at \m{y\approx x}, for a combination of the
centred sums of §2. It is not yet $(\star)$: the combination involves $\chid$ instead of~$\psi$, and with a minus
sign. Unequal weights repair this, in §5.

In the archive the outcome of this section is a single proposition, {\leanfont TopG}: whenever hosts exist on
arbitrarily late windows and satisfy the cohomological bound of §6, they carry chains for which the inequality
holds, with the exponents hidden in $x^{O(\eps)}$ fixed before~$\eps$ is chosen.
\leanline{hosts \lean{ischifield}\sep integrality \lean{chainint}\sep local divisibility \lean{minor}\sep slopes
\lean{slope}\sep periods \lean{chowla}\sep the two sides \lean{ask}\sep sizes \lean{rank}\sep the outcome \lean{topg}}

% ------------------------------------------------------------------------------------- 5
\sect{5}{est}{The estimate, II: from the slope inequality to $(\star)$}
\textbf{The characters of the host.} The extension $E/\Q$ is abelian, and
\[
\Gal(E/\Q)\cong(\Z/2)^2\times A,\qquad A=\Gal\bigl(\Q(\zeta_{2^{e+2}})^+K_\psi/\Q\bigr),
\]
the two factors $\Z/2$ being the Galois groups of $\Q(\sqrt{-q})$ and of~$K$. The characters of $\Gal(E/\Q)$ are
Dirichlet characters: the products $\chiq^{\,i}\,\chid^{\,j}\,\omega$ with \m{i,j\in\{0,1\}} and $\omega$ a character
of~$A$. Let \m{A_0\subset A} be the subgroup that fixes~$K_\psi$. The characters of~$A$ that are trivial on~$A_0$ are
exactly the $o$ powers of~$\psi$.

\textbf{What unequal weights do.} For general weights the difference $\mathrm{ASK}-\mathrm{MAIN}$ is a finite Fourier
sum of centred sums. Write $w_\sigma$ for the weight attached to an element~$\sigma$ of $\Gal(E/\Q)$ outside~$\Phi$.
Then
\[
\mathrm{ASK}(a)-\mathrm{MAIN}(a)=\frac1{\lvert\Phi\rvert}\ \sum_{\chi}\ \sum_{\sigma\notin\Phi}\,
\overline{\chi(\sigma)}\;w_\sigma\;g_\chi\Bigl(\frac a{w_\sigma}\Bigr),
\]
where $\chi$ runs over all characters of $\Gal(E/\Q)$, with \m{g_1=G}. For equal weights only \m{\chi=1} and
\m{\chi=\chid} survive, and this is the formula of §4. In general the weights decide with which coefficient each
character occurs.

The construction of §4 comes as a family. Its members are indexed by the elements~$\tau$ of a finite set
\m{T\subset A}, called tops in the archive, and by translations by the elements of~$A_0$. Passing from one member to
another moves the weights along the group, and so multiplies the coefficient of each character~$\chi$ by a value
of~$\chi$. Three steps turn the family of inequalities into~$(\star)$.
\begin{enumerate}[label=\textup{(\roman*)}]
\item \kw{Averaging over $A_0$.} Adding the inequalities over all translations by~$A_0$ annihilates every character
that is non-trivial on~$A_0$, by the orthogonality of characters. There remain $4o$ characters: the powers of~$\psi$
and their products with $\chiq$, $\chid$ and $\chiq\chid$.
\item \kw{Positivity by design.} The $o$ powers of~$\psi$ are the characters wanted: the primitive characters
inducing them form~$P(\psi)$. After averaging over \m{\tau\in T}, a power $\psi'$ of~$\psi$ carries the coefficient
\m{\lvert T\rvert^{-1}\sum_{\tau\in T}\psi'(\tau)}. The set~$T$ is chosen so that the proportion of its elements lying
in~$A_0$ is $4/(3o+1)$, and the proportion in each of the other \m{o-1} cosets of~$A_0$ is $3/(3o+1)$, up to a
negligible error. The coefficient is then $1$ for \m{\psi'=1} and $1/(3o+1)$ for \m{\psi'\ne1}. These are the positive
weights of $P(\psi)$ in §3.
\item \kw{Twisting.} The other $3o$ characters occur with coefficients of no definite sign. They are removed one at a
time by means of the bound \m{\lvert g_\chi(x)\rvert\le x^{1/2+\eps}}. For a given character such a bound is as hard
as the theorem itself; but the primes $q$ and~$d$ are ours to choose. An averaging argument over the candidate primes
(a mean-square bound) shows that for every window some pair $(q,d)$ satisfies the bound for all $3o$ twisted
characters, at every point of the window.
\end{enumerate}

\textbf{The auxiliary primes.} For a window beginning at~$Y$ put \m{z=Y^{\eps/5}}. The primes $q$ and~$d$ are
required to satisfy \m{\frac z2<\log q\le z}, \m{\frac z2(\log Y)^2<\log d\le z(\log Y)^2},
\[
q\equiv3\pmod 8,\qquad d\equiv7\pmod{16},\qquad q\mid d+1,
\]
and every odd prime \m{p\le z} must divide \m{d+1}. These conditions make $2$, $q$ and every odd prime up to~$z$
split in~$K$, so none of them occurs in $\mathrm{ASK}$. As the pair $(q,d)$ depends on the window, so does the host~$E$: hosts are rebuilt at every scale.

\textbf{The thick set.} The family involves one more finite set, a subset $T_2$ of~$A_0$ that the archive calls
thick. It must contain at least the proportion $x^{-\eps}$ of~$A$ and be free of short linear relations. Such sets come
from Behrend's construction, which uses the integer points on a sphere: for integers \m{Q\ge1} and \m{H\ge2} there
are at least \m{Q\exp(-\kappa\sqrt{\log Q\,\log H})/H} non-negative integers below $Q/(4H)$, with an absolute
constant~$\kappa$, no three distinct ones $s_1$, $s_2$, $s_3$ of which satisfy
\[
j\,(s_2-s_1)=i\,(s_3-s_1)\qquad\text{with }i\ne j\text{ and }0<\lvert i\rvert,\lvert j\rvert\le H.
\]
For \m{(i,j)=(1,2)} this excludes three-term arithmetic progressions. It is applied with \m{Q=\lvert A_0\rvert}
and $H$ of the size of $\log x$. As a consequence, the orbit of a point under a Frobenius element meets a translate
of~$T_2$ at most twice in any $H$ consecutive steps.

\textbf{Bookkeeping.} In units of~$\eps$, the exponent lost is $\frac{13}8$ in the sizes of~$m$ and of the family, $5$
in the error terms of the slope inequality, $1$ in the density of the thick set, $\frac38$ in constants, and $1$ in
removing the twisted characters. The sum is \m{\frac{13}{8}+5+1+\frac38+1=9}, the $9$ of~$(\star)$.
\leanline{Fourier identity \lean{chbalance}\sep weights \lean{shares}\sep twisted characters \lean{bounded}\sep
auxiliary primes \lean{pairex}\sep hosts at every scale \lean{supply}\sep the thick set \lean{behrend}\sep conclusion
\lean{core-gm}}

% ------------------------------------------------------------------------------------- 6
\sect{6}{arith}{The price of divisibility: Iwasawa theory}
The divisibility of §4(a) holds only up to a loss at each lift, and the loss is governed by Galois cohomology. For a
prime~$\ell$ and an integer~$h$, let $\Q_\ell/\Z_\ell(-h)$ be the group $\Q_\ell/\Z_\ell$ on which Galois groups act
through the $(-h)$-th power of the $\ell$-adic cyclotomic character, and let $G_{E,\ell}$ be the Galois group of the
maximal extension of~$E$ unramified outside~$\ell$. Put
\[
\Sha^1_\ell(E,h)=\ker\Bigl(H^1\bigl(G_{E,\ell},\Q_\ell/\Z_\ell(-h)\bigr)\longrightarrow
\prod_{v\mid\ell}H^1\bigl(E_v,\Q_\ell/\Z_\ell(-h)\bigr)\Bigr),
\]
the group of classes that are trivial at every place~$v$ of~$E$ above~$\ell$, where $E_v$ is the completion of~$E$
at~$v$. Write $\len$ for the length of a $\Z_\ell$-module. By Poitou--Tate duality $\Sha^1_\ell(E,h)$ has the same
length as the analogous group $\Sha^2$ with coefficients $\Z_\ell(h+1)$.

Vandiver's conjecture, that $\ell$ does not divide the class number of $\Q(\zeta_\ell)^+$, is a statement on the
vanishing of groups of this kind; the announcement speaks of ``a significant weakening of Vandiver''. What is used is
not a vanishing but a bound for the total length, on a whole window:

\statement{arith}{\textbf{Statement} ({\leanfont VWinAt}\,$\psi$\,{\leanfont 401}).\enspace\emph{For every
$0<\eps<\frac1{52}$ there is $J\ge1$ such that for all sufficiently large~$Z$, and every pair of primes $(q,d)$
satisfying the conditions of §5 with $Z^{1/J}$ in place of~$Y$, there is $Y$ with $Z^{1/J}\le Y\le Z^{1/2}$ such that
for every $x\in[Y,Y^2]$ the host~$E$ at scale~$x$ satisfies}\par\vspace{5pt}
\centerline{$\displaystyle\sum_{\ell}\,\log\ell\!\!\sum_{1\le h<x^{401}}\!\!\len\,\Sha^1_\ell(E,h)\ \le\
\exp\bigl(x^{\eps}/2\bigr)$,}\par\vspace{4pt}
\emph{the outer sum running over all primes~$\ell$.}}

\textbf{Why such a window exists.} Fix~$\ell$. Along a cyclotomic tower these lengths are the lengths of the
cokernels of an operator, twisted by~$h$, on an Iwasawa module; by the theorem of Ferrero and Washington the module
is finitely generated over~$\Z_\ell$. Let $R$ bound its rank. Then there are at most $R$ numbers \m{c_1,\dots,c_R} in
$\Z_\ell$, the \kw{centres}, with
\[
\len\,\Sha^1_\ell(E,h)\ \le\ \sum_{i=1}^{R}v_\ell(c_i-h)\ +\ (\text{terms depending only on the layer}),
\]
where $v_\ell$ is the $\ell$-adic valuation. So a twist~$h$ is expensive only if it is $\ell$-adically very close to
a centre. For a bound \m{\ell^{k}\le Z^{401}<\ell^{k+1}}, call an integer \m{1\le h\le Z^{401}} \kw{exceptional} if
\m{v_\ell(c_i-h)>k} for some~$i$. Two exceptional integers at the same centre differ by a multiple of~$\ell^{k+1}$,
which exceeds $Z^{401}$, so they are equal. Each centre therefore accounts for at most one exceptional twist: there
are at most~$R$ of them, whatever~$Z$.

An exceptional twist~$h$ does no harm at a scale~$x$ where it is out of range, \m{h>x^{401}}, or small,
\m{h\le x^{\nu}} for a fixed small \m{\nu>0}. In the variable $\log x$ it therefore forbids one interval
$[\lambda,B\lambda)$, with \m{\lambda=\frac1{401}\log h} and \m{B=401/\nu}.

\statement{arith}{\textbf{Interval lemma.}\enspace\emph{Let $B\ge1$ and $V>0$. For any $R$ intervals
$[\lambda_i,B\lambda_i)$ there is an interval $[V',2V']$ disjoint from all of them, with}\quad
$\dfrac{V}{2^{R+2}B^{R}}\le V'$\quad\emph{and}\quad$2V'\le V$.}

With \m{V=\log Z} and \m{Y=\exp(V')}, the interval $[Y,Y^2]$ contains no forbidden scale and
\m{Z^{1/J}\le Y\le Z^{1/2}} for \m{J=4\,(2B)^R}. This is the window.
\leanline{the window \lean{vwin-def}\sep interval lemma \lean{window1}}

% ------------------------------------------------------------------------------------- 7
\sect{7}{arith}{The automorphic foundation}
Both branches of the diagram start from Galois representations attached to automorphic forms: the cohomological bound
of §6 is deduced from them, by base change, and so are inputs of §4. The archive constructs them in the following
generality. Let $E$ be any CM field, \m{n\ge1}, $\mathbb A_E$ the ring of adèles of~$E$, and $\pi$ a cuspidal
automorphic representation of $\mathrm{GL}_n(\mathbb A_E)$ that is
\begin{itemize}[label={\color{arith}\textbullet}]
\item \kw{regular algebraic}: at each infinite place its infinitesimal character is that of an irreducible algebraic
representation of~$\mathrm{GL}_n$, so that its exponents are distinct and lie in \m{\frac{n-1}2+\Z};
\item \kw{polarized}: it is conjugate self-dual up to a twist, \m{\pi^c\cong\pi^\vee\otimes(\chi\circ N_{E/E^+})},
for a Hecke character~$\chi$ of~$E^+$ of the appropriate parity, where $N_{E/E^+}$ is the norm.
\end{itemize}
Fix a prime~$\ell$ and an isomorphism \m{\iota\colon\overline{\Q}_\ell\cong\mathbb C}. Then there is a continuous
semisimple representation
\[
r_\pi\colon\Gal(\overline E/E)\longrightarrow\mathrm{GL}_n(\overline{\Q}_\ell)
\]
which, at all but finitely many places~$v$ of~$E$, is unramified and satisfies
\[
\iota\,\operatorname{charpoly}\,r_\pi(\mathrm{Frob}_v^{-1})=\prod_{\alpha}\bigl(X-q_v^{(n-1)/2}\alpha\bigr).
\]
Here $\mathrm{Frob}_v$ is a Frobenius element at~$v$, $q_v$ is the number of elements of the residue field, and
$\alpha$ runs over the Satake parameters of the local component~$\pi_v$, the $n$ complex numbers that determine an
unramified~$\pi_v$. In addition, $r_\pi$ corresponds to~$\pi_v$ under the local Langlands correspondence at the
places \m{v\nmid\ell}, in the cases specified there; and at the places \m{v\mid\ell} it is de~Rham, and crystalline
where $\pi_v$ is unramified, with Hodge--Tate weights given by the infinitesimal character of~$\pi$ (these are the
conditions of $p$-adic Hodge theory satisfied by representations that come from geometry). These are the
properties known from the work of Harris--Taylor, Shin, Chenevier--Harris, Caraiani and others; the archive contains
its own formal proof of the statement it uses.

The construction starts from two inputs, a geometric one resting on a holomorphic Lefschetz fixed-point formula, and
a theorem on theta lifts. It proceeds by congruences and limits of traces (pseudocharacters). It first treats CM
fields containing an imaginary quadratic field, under further conditions, and then removes these: $\pi$ is twisted
to adjust the polarization and transferred by base change to infinitely many quadratic CM extensions
$E(\sqrt{-p})$, and the representations obtained there are patched into~$r_\pi$ by comparing traces.

The final assembly of the theorem has seven groups of hypotheses, the ``seven inputs'': facts on the unitary group at
an auxiliary place, local theta correspondences, the agreement of theta integrals, a local matching statement, a
covolume bound, the existence of local Langlands parameters, and the cohomological window of §6.
\leanline{the Galois representations \lean{row-def}, \lean{row-prod}\sep the final assembly \lean{seven}}

% ------------------------------------------------------------------------------------- synthesis
\head{Why each idea calls for the next}
\begin{list}{}{\setlength{\leftmargin}{3.3em}\setlength{\labelwidth}{2.75em}\setlength{\labelsep}{0.55em}%
  \setlength{\itemsep}{5pt}\setlength{\parsep}{0pt}\setlength{\topsep}{3pt}\setlength{\partopsep}{0pt}}
\item[\dotnum{2}{ana}\hspace{1.5pt}\dotnum{3}{ana}\hfill] A proof by upper bounds needs a statistic whose transform
sees every zero with one sign: hence ramps, centring and positive families.
\item[\dotnum{4}{est}\hfill] Geometry yields upper bounds through slopes, and slopes see primes through divisibility.
Divisibility in bulk needs multiplicity, which iterated theta lifts supply at a price that fixes their number at
$\sqrt x$, and whose archimedean side reproduces the centring.
\item[\dotnum{5}{est}\hfill] Theta lifts live on unitary groups of CM fields, so the bound arrives for a whole Galois
group of characters at once: hence averaging, positivity by design, and twisting.
\item[\dotnum{6}{arith}\hspace{1.5pt}\dotnum{7}{arith}\hfill] Divisibility has a cohomological defect, tamed by Iwasawa
theory in the very tower that carries the averaging; both rest on the Galois representations of the forms involved.
\item[\dotnum{3}{ana}\hfill] Such a construction exists only at special scales, so the analytic theorem must be local
in~$y$.
\end{list}

% ------------------------------------------------------------------------------------- status
\par\addvspace{12pt plus 4pt}
\Needspace{10\baselineskip}
\begingroup
\renewcommand{\leanfont}{\ttfamily\fontsize{7.9}{10}\selectfont\color{leancol}}
\setlength{\fboxsep}{9pt}\setlength{\fboxrule}{0.6pt}
\noindent\fcolorbox{rulegrey}{ink!2!white}{\parbox{\dimexpr\linewidth-2\fboxsep-2\fboxrule\relax}{%
\fontsize{9.4}{12.8}\selectfont
{\bfseries What is machine-checked}\quad{\color{inksoft}records of \AsOf}\par\vspace{3pt}
{\color{inksoft}Lean is a proof assistant: a proof is a formal object, and a small program, the kernel, checks it.
Mathlib is Lean's library of mathematics; {\leanfont RiemannHypothesis} is its statement of RH. NanoDA is an
independent kernel, written in Rust.\par}\vspace{4pt}
\begin{list}{}{\setlength{\leftmargin}{2.1em}\setlength{\labelwidth}{1.6em}\setlength{\labelsep}{0.5em}%
  \setlength{\itemsep}{3pt}\setlength{\parsep}{0pt}\setlength{\topsep}{0pt}\setlength{\partopsep}{0pt}}
\item[\textcolor{ok}{\checkmark\checkmark}\hfill] $(\star)$ for $\psi=1$ $\Rightarrow$ {\leanfont RiemannHypothesis}:
Lean kernel and NanoDA (\NanodaDecls{} declarations).
\item[\textcolor{ok}{\checkmark}\hfill] {\leanfont TopG} (§4) and {\leanfont VWinAt}\,$\psi$\,{\leanfont 401} for
every~$\psi$ (§6) $\Rightarrow$ Theorem: Lean kernel (\DirichletModules{} modules).
\item[\FullProofIcon\hfill] The five theorems themselves, in the Lean archive of L.\,Alpöge (with Claude):
\FullProofLine
\item[\textcolor{inksoft}{$\circ$}\hfill] §§4--7 are read here from the Lean source and from the announcement, not
re-derived.
\item[\textcolor{ok}{\checkmark}\hfill] Checked in this repository: the logical skeleton of the guide
({\leanfont lean/Atlas/Architecture.lean}, no axioms) and §1 for $\zeta$ against Mathlib at the archive's pin
({\leanfont lean-mathlib/AtlasZeta.lean}).
\end{list}}}
\endgroup

% ------------------------------------------------------------------------------------- notation
\newpage
\head{Notation}
{\fontsize{9.6}{12.6}\selectfont\color{inksoft}\noindent The symbols of the text, in the order of their first
appearance. Bound variables of sums and products ($n$, $p$, $i$, $j$) are omitted; a letter in brackets is the one
used in the archive or in the announcement.\par}
\vspace{4pt}
\begingroup\fontsize{8.6}{10.9}\selectfont\raggedright
\renewcommand{\leanfont}{\ttfamily\fontsize{7.6}{10}\selectfont\color{leancol}}
\setlength{\columnsep}{9mm}
\begin{multicols}{2}
\nt{$s$, $y$}{complex variable, $\Re s$ its real part; real variable \m{y\ge1}}
\nt{$O(\cdot)$, $\asymp$, $\approx$}{at most a constant times; of the same order; close to}
\nt{$\F$, $X$}{a finite field; in the announcement's sentence on page~1, the length of the prime sum}
\nt{$\chi$, $\bar\chi$, $\chi^*$}{a Dirichlet character, its complex conjugate, the primitive character inducing it; $N$ its modulus}
\nt{$L(s,\chi)$, $\zeta(s)$}{Dirichlet $L$-function; the case $N=1$}
\nt{$\Lambda(s,\chi)$, $W(\chi)$}{completed $L$-function, root number (§1)}
\nt{$\psi$, $o$}{the primitive character under study; its order}
\nt{$\Lambda(n)$, $\gamma$}{von Mangoldt function; Euler's constant (§2)}
\nt{$(u)_+$}{$\max(u,0)$}
\nt{$G(y)$, $g_\chi(y)$}{centred ramp sums of $\zeta$ and of~$\chi$ (§2)}
\nt{$\rho$, $m_\chi(\rho)$}{a zero; its multiplicity as a zero of $L(s,\chi)$}
\nt{$F$, $c_\chi$, $G_F$}{positive family, its weights, its sum (§3) [weights $a$ in the archive]}
\nt{$\theta$, $\delta$, $Y$, $C$}{exponent under test; width and left end of a window $[Y,Y^{1+\delta}]$; a constant}
\nt{$S(\rho)$}{weighted multiplicity of~$\rho$ in a family (§3)}
\nt{$g(u)$, $\widehat g$, $A_k$, $\beta$}{$G_F(e^u)$; its Laplace transform; its normalized $k$-th derivative; $\Re s$ (§3)}
\nt{$\eta$, $C_\eta$}{an arbitrarily small exponent and its constant (§3)}
\nt{$P(\psi)$, $G_\psi$}{the family of the powers of~$\psi$; its sum (§3)}
\nt{$\eps$, $\eps_1$}{small exponent, chosen after~$\theta$; its upper limit}
\nt{$x$}{the scale: the point at which $(\star)$ is proved}
\nt{$(\star)$}{the arithmetic statement (§3)}
\nt{$\vartheta$, $D$, $\Omega$}{theta function, invariant derivation, period of a CM elliptic curve (§4)}
\nt{$a$}{height: order of derivative, then the parameter of $\mathrm{ASK}$ and $\mathrm{MAIN}$}
\nt{$q$, $d$, $z$}{auxiliary primes (§4, §5); $z=Y^{\eps/5}$, the size of $\log q$}
\nt{$K$, $\Phi$, $\chid$, $\chiq$}{$\Q(\sqrt{-d})$; $\Gal(E/K)$; the quadratic characters of $K$ and of $\Q(\sqrt{-q})$}
\nt{$e$}{integer fixing the layer: $2^e\asymp x^{2+8\eps}$ (in $e^u$ it is the usual constant)}
\nt{$E$, $E^+$, $c$}{host field, its maximal real subfield, complex conjugation (§4)}
\nt{$\zeta_n$, $K_\psi$}{primitive $n$-th root of unity; field cut out by~$\psi$}
\nt{$X_j$, $U(X_j)$}{Hermitian spaces over~$E$ and their unitary groups (§4)}
\nt{$n_0$, $t$, $m$, $\delta_1$}{$\dim X_0$ [$k$ in the archive]; odd integer with $\dim X_{j+1}-\dim X_j=t+1$; number of lifts; $x^{-\eps/8}$}
\nt{$M$, $f$}{class of Hermitian lattices; automorphic form}
\nt{$\xi_j$, $\vartheta_j$, $\Theta_j$, $\Theta$}{Hecke characters [$\mu_j$], theta kernels, theta lifts, the chain}
\nt{$\pi$, $\pi^c$, $\pi^\vee$, $\boxplus$}{automorphic representation, its conjugate, its contragredient; isobaric sum}
\nt{$\varpi$}{an element with $\varpi^2=p$ (§4a)}
\nt{$\mathrm{ASK}$, $\mathrm{MAIN}$}{the finite and the archimedean side of the slope inequality (names of the archive)}
\nt{$w$, $w_n$, $w_\sigma$, $b$}{weights; their common size $\exp(x^{\eps})$}
\nt{$\mathcal L$, $\mu$, $\mu_{\max}$, $\lVert\varphi\rVert$}{Euclidean lattice, slope, maximal slope, operator norm (§4b)}
\nt{$A$, $A_0$, $\omega$, $\sigma$}{$\Gal(\Q(\zeta_{2^{e+2}})^+K_\psi/\Q)$; its subgroup fixing~$K_\psi$; a character of~$A$; an element of $\Gal(E/\Q)$ (§5)}
\nt{$T$, $\tau$, $T_2$}{the set of tops, one of its elements; the thick set (§5)}
\nt{$Q$, $H$, $\kappa$}{size and relation bound in Behrend's construction, its constant (§5)}
\nt{$\ell$, $h$, $v$}{a prime; a Tate twist [$e$ in the archive]; a place of~$E$}
\nt{$G_{E,\ell}$, $E_v$, $H^1$}{Galois group unramified outside~$\ell$; completion; Galois cohomology}
\nt{$\Sha^1_\ell(E,h)$, $\len$}{locally trivial classes (§6); length over $\Z_\ell$}
\nt{$Z$, $J$}{large parameter; exponent with $Z^{1/J}\le Y$ (§6)}
\nt{$c_i$, $R$, $v_\ell$}{centres, a bound for their number, $\ell$-adic valuation}
\nt{$\nu$, $B$, $\lambda_i$, $V$, $V'$}{quantities of the interval lemma (§6)}
\nt{$n$, $\mathbb A_E$}{rank; adèles of~$E$ (§7)}
\nt{$r_\pi$, $\iota$}{Galois representation of~$\pi$; fixed isomorphism $\overline{\Q}_\ell\cong\mathbb C$}
\nt{$\mathrm{Frob}_v$, $q_v$, $\alpha$}{Frobenius element, size of the residue field, Satake parameters}
\nt{{\leanfont TopG}}{the outcome of §4 as one proposition of the archive}
\nt{{\leanfont VWinAt}\,$\psi$\,{\leanfont 401}}{the cohomological window of §6}
\end{multicols}
\endgroup

\end{document}
